\chapter{Coughs and Colds (in Children)}

\section*{Definition and Overview}
Coughs and colds are the most common illnesses of childhood. A cold is a viral infection of the nose and throat, and a cough is the body's way of clearing the airways. Children have more colds than adults -- typically six to ten a year -- because their immune systems are still developing and they are exposed to many viruses at school, daycare, and home. Most coughs and colds in children are mild and settle within one to two weeks without specific treatment. The main tasks for parents are to keep the child comfortable, watch for complications, and know when to seek medical help.

\section*{Epidemiology}
Viral respiratory infections are the most frequent reason children see a doctor. Children under six typically have six to eight colds per year, with the highest rates in those attending childcare. Colds are more common in autumn and winter, and are spread through coughs, sneezes, and contaminated hands. Most are caused by rhinoviruses and other respiratory viruses. Although the illnesses themselves are usually mild, they are a major cause of school and work absence, and can be more serious in babies, children with chronic conditions, and those born prematurely.

\section*{Aetiology and Pathophysiology}
Many different viruses cause colds and coughs, and they infect the lining of the nose, throat, and airways.

\begin{itemize}
    \item \textbf{Rhinoviruses:} the most common cause of the common cold
    \item \textbf{Other respiratory viruses:} adenovirus, respiratory syncytial virus (RSV), influenza, parainfluenza, and coronaviruses
    \item \textbf{Transmission:} droplets from coughing and sneezing, and hand-to-face contact
    \item \textbf{Infection:} the virus multiplies in the airway lining, causing inflammation, mucus, and coughing
    \item \textbf{Recovery:} the immune system clears the virus over a week or two
    \item \textbf{Secondary infection:} viruses can damage the airway lining, allowing bacterial infection to follow
\end{itemize}

\section*{Clinical Features}
Symptoms develop over one to three days and typically last one to two weeks.

\begin{itemize}
    \item Runny or blocked nose and sneezing
    \item Sore throat and mild fever
    \item Cough, which may be dry at first and become productive
    \item Reduced appetite and irritability
    \item Mild tiredness or generally feeling unwell
    \item Watery eyes, mild headache, or ear discomfort
    \item A cough that can last for two to three weeks after the other symptoms settle
\end{itemize}

\section*{Differential Diagnosis}
Other childhood conditions can present with cough, fever, and coryza.

\begin{itemize}
    \item Influenza, which tends to be more severe with higher fever and body aches
    \item Bronchiolitis, especially in babies under one year, with fast or laboured breathing
    \item Pneumonia, with high fever, fast breathing, and chest signs
    \item Croup, with a barking cough and noisy breathing
    \item Asthma, with wheeze, cough at night, or cough with exercise
    \item Whooping cough (pertussis), with bouts of coughing and a whoop
    \item Ear infections, tonsillitis, and sinusitis
    \item A foreign body inhaled into the airway (in a child with sudden coughing)
\end{itemize}

\section*{When to Seek Medical Care}
Most coughs and colds are managed at home, but medical review is needed if the child is very young, is not drinking, or seems more than mildly unwell.

\textbf{Seek medical advice if:}
\begin{itemize}
    \item The child is under three months old and has a fever or is unwell
    \item The fever lasts more than a few days or is very high
    \item The child is breathing fast, working hard to breathe, or has noisy breathing
    \item The child is not drinking enough, is passing little urine, or shows signs of dehydration
    \item Symptoms worsen after several days, or the child seems very unwell
    \item There is ear pain, a persistent cough, or the cough lasts more than three weeks
\end{itemize}

\textbf{Call 000 (or local emergency services) for:}
\begin{itemize}
    \item Difficulty breathing, blue lips, or pauses in breathing
    \item A child who is floppy, unusually drowsy, or difficult to wake
    \item A seizure
    \item A rash that does not fade when pressed
    \item Choking or suspected inhalation of an object
\end{itemize}

\section*{Diagnosis and Investigations}
Most coughs and colds are diagnosed clinically and need no tests.

\begin{itemize}
    \item \textbf{History and examination:} the pattern of symptoms, fever, feeding, and breathing
    \item \textbf{Oxygen monitoring:} pulse oximetry if breathing is a concern
    \item \textbf{Chest X-ray:} only if pneumonia or another complication is suspected
    \item \textbf{Swabs and blood tests:} for influenza, RSV, or whooping cough in specific situations
    \item \textbf{Review:} of vaccination status, particularly for whooping cough and influenza
\end{itemize}

\section*{Management}
Antibiotics do not help viral colds, and treatment is supportive.

\begin{itemize}
    \item \textbf{Fluids:} encourage regular drinks to prevent dehydration
    \item \textbf{Fever relief:} paracetamol or ibuprofen, using the correct dose for age and weight
    \item \textbf{Comfort:} rest, and clearing a blocked nose with saline drops or spray
    \item \textbf{Honey:} for children over one year, a small amount can soothe a cough
    \item \textbf{Humid air:} a cool-mist humidifier or steam in the bathroom may ease cough
    \item \textbf{Over-the-counter cold medicines:} not recommended for children under six, and used with care in older children
    \item \textbf{Avoid smoke:} keep the child away from cigarette and vape smoke
    \item \textbf{Treat complications:} antibiotics for bacterial infections such as ear infections when diagnosed
\end{itemize}

\section*{Complications}
\begin{itemize}
    \item Middle ear infection, the most common complication
    \item Sinusitis and tonsillitis
    \item Pneumonia or worsening of asthma
    \item Bronchiolitis and croup in babies
    \item Dehydration from poor feeding and fever
    \item Febrile seizures in some children
    \item Spread of the infection to other family members
\end{itemize}

\section*{Prognosis and Outcomes}
Almost all children recover fully from coughs and colds within one to two weeks, though a mild cough can persist for up to three weeks. Repeated colds through childhood are normal and do not indicate a weak immune system. Complications such as ear infections are usually treated easily. Antibiotics are not needed for uncomplicated colds, and their overuse contributes to resistance. Most children grow out of the high frequency of colds as they reach school age.

\section*{Prevention and Self-Care}
\begin{itemize}
    \item Encourage regular handwashing with soap and water
    \item Teach children to cough or sneeze into a tissue or their elbow
    \item Keep children home from childcare or school while significantly unwell
    \item Keep vaccinations up to date, including influenza and whooping cough
    \item Keep babies and children away from cigarette and vape smoke
    \item Maintain good nutrition, sleep, and physical activity
    \item Treat fever and discomfort with age-appropriate doses of paracetamol or ibuprofen
    \item Know the warning signs that need medical review
\end{itemize}

\section*{Sources and Further Reading}
The following reputable sources provide further information for healthcare students and patients seeking reliable, current advice about coughs and colds in children.

\begin{itemize}
    \item Healthdirect Australia. \textit{Coughs and colds in children.} https://www.healthdirect.gov.au/
    \item Raising Children Network. \textit{Colds and coughs in children.} https://raisingchildren.net.au/
    \item NHS. \textit{Coughs and colds in children.} https://www.nhs.uk/conditions/coughs-and-colds-in-children/
    \item Royal Children's Hospital Melbourne. \textit{Kids Health Info: colds and coughs.} https://www.rch.org.au/kidsinfo/
    \item Centers for Disease Control and Prevention. \textit{Common cold in children.} https://www.cdc.gov/
    \item Sydney Children's Hospitals Network. \textit{Childhood respiratory infections.} https://www.schn.health.nsw.gov.au/
\end{itemize}
